\documentclass[12pt,a4paper]{article}
\usepackage[spanish,es-nodecimaldot]{babel}
\usepackage[utf8]{inputenc}
\usepackage[T1]{fontenc}
\usepackage{geometry}
\usepackage{hyperref}
\usepackage{listings}
\usepackage{xcolor}
\geometry{margin=2.5cm}
\definecolor{codegray}{gray}{0.95}
\lstset{language=R,basicstyle=\ttfamily\small,backgroundcolor=\color{codegray},frame=single,breaklines=true}
\title{Almacén de datos (Data Warehouse)}
\author{Unidad 1}
\date{2026}
\begin{document}
\maketitle
\section{Definición}
Un Data Warehouse es un repositorio integrado, histórico y orientado al análisis. Reúne datos de distintas fuentes para facilitar informes, inteligencia de negocios y minería de datos.
\section{Componentes}
Sus componentes habituales son las fuentes operacionales, los procesos ETL (extraer, transformar y cargar), el almacenamiento central, los metadatos y las herramientas de consulta y visualización.
\section{Relación con minería de datos}
El almacén de datos proporciona datos consistentes y consolidados para explorar patrones, construir modelos predictivos y comparar indicadores empresariales a lo largo del tiempo.
\section{Modelo dimensional}
Un diseño frecuente separa tablas de hechos, que contienen medidas como cantidad e ingreso, y tablas de dimensiones, que describen productos, fechas, clientes o regiones. El esquema estrella facilita consultas analíticas.
\section{ETL y calidad}
El proceso ETL debe validar tipos, eliminar duplicados, controlar valores faltantes y conservar la trazabilidad desde la fuente hasta la tabla final. Las reglas de calidad deben quedar documentadas y monitorearse de forma periódica.
\section{Data Warehouse y Data Lake}
Un Data Warehouse prioriza datos estructurados, consistentes y listos para consulta. Un Data Lake puede conservar datos en formatos diversos con mayor flexibilidad. Ambos pueden integrarse en una arquitectura moderna, según las necesidades de gobierno y análisis.
\section{Ejemplo en R}
\begin{lstlisting}
# Simulacion de una tabla dimensional para ventas
ventas <- data.frame(
  fecha = as.Date(c("2026-01-01", "2026-01-02", "2026-01-02")),
  producto = c("A", "A", "B"),
  cantidad = c(3, 2, 5),
  ingreso = c(30, 20, 75)
)
aggregate(cbind(cantidad, ingreso) ~ producto, data = ventas, sum)
\end{lstlisting}
\section{Referencias iniciales}
\begin{itemize}
\item Nambiar, A. y Mundra, D. (2022). An Overview of Data Warehouse and Data Lake in Modern Enterprise Data Management. DOI: \url{https://doi.org/10.3390/bdcc6040132}.
\item Deb Nath, R. P. y Hose, K. (2017). SETL: A programmable semantic extract-transform-load framework. DOI: \url{https://doi.org/10.1016/j.is.2017.01.005}.
\end{itemize}
\end{document}
